% =====================================================================
% 第 I 章  Statement of results (J-P. Serre, Oct. 16, 1957)
% =====================================================================
\lecture{Statement of results}{Statement of Results}{J-P. Serre}{Oct. 16, 1957}

% ---------------------------------------------------------------
\section{The notion of complex multiplication}

Let $X$ be an elliptic curve. As a complex Lie group, it is the quotient
of the complex plane $\mathbb{C}$ by a lattice $\Gamma$, spanned by two
periods $\omega_{1}$, $\omega_{2}$, and since $X$ is isomorphic to the
curve defined by the periods $z\omega_{1}$, $z\omega_{2}$ for any non-zero
$z\in\mathbb{C}$, we may assume $\Gamma$ to be spanned by $1$ and $\tau$,
where $\tau$ has a positive imaginary part.

An endomorphism of $X$ may be identified with an endomorphism of its
universal covering $\mathbb{C}$ mapping $\Gamma$ into itself; it is
therefore the multiplication by a complex number $z$ such that
$z,\ z\tau\in\Gamma$. The endomorphisms of $X$ form a ring $A(X)$, which
always contains the integers $\mathbb{Z}$ (the ``trivial endomorphisms'').
The other ones (if any) are given by complex numbers and are called
\emph{complex multiplications}. If $A(X)\neq\mathbb{Z}$, the curve $X$ is
said to admit complex multiplications.

``In general'', $X$ has no complex multiplication. In fact, assume that
$z$ defines a non-trivial endomorphism of $X$. Then
\[
z=a+b\tau,\qquad z\tau=c+d\tau,\qquad(a,b,c,d\ \text{integers},\ b\neq 0),
\]
whence
\[
a\tau+b\tau^{2}=c+d\tau
\]
and $\tau$ must belong to an imaginary quadratic field, say $K$; moreover
$z$ belongs to the ring of integers $\mathfrak{o}(K)$ of $K$, since it is
in $K$ and defines an endomorphism of a $\mathbb{Z}$-module of finite
rank, namely $\Gamma$. Therefore, $A(X)$ is an \emph{order} of $K$
(subring of $\mathfrak{o}(K)$ containing $\mathbb{Z}$ and which has rank
$2$ as a $\mathbb{Z}$-module);

one gets in this way all orders of all quadratic imaginary fields (if $R$
is such an order, take for $X$ a curve with lattice of periods $R$; since
$1\in R$, $zR\subset R$ if and only if $z\in R$, whence $A(X)=R$).

Assume that $A(X)=\mathfrak{o}(K)$ and that $\Gamma\subset K$. Then
$\Gamma$ is an ideal of $K$, and conversely any ideal of $K$ gives rise
to a curve $X$ such that $A(X)=\mathfrak{o}(K)$. Two such curves are
isomorphic if and only if the corresponding ideals are homothetic,
i.e.\ belong to the same \emph{ideal class}.

Let $j$ be the modular function. For the curve with normal equation
\[
y^{2}=4x^{3}-g_{2}x-g_{3}
\]
its value is
\[
j=2^{6}\cdot 3^{3}\cdot g_{2}^{3}/\Delta,\qquad
(\Delta=g_{2}^{3}-27\cdot g_{3}^{2}\,).
\]

Two elliptic curves are isomorphic over an algebraically closed field if
and only if their modular invariants are equal. By the above, $j$ defines
a function on the ideal classes $\mathfrak{k}_{1},\ldots,\mathfrak{k}_{n}$
of $K$; the numbers $j(\mathfrak{k}_{i})$ are ``singular values'' of $j$,
and are called the \emph{class invariants} of $K$; they are pairwise
different, and have proved to be of fundamental importance in the study
of the abelian extensions of $K$, to which we now turn.

% ---------------------------------------------------------------
\section{Unramified abelian extensions of an imaginary quadratic field}

It is a classical result of Kronecker that every abelian extension
(i.e.\ normal extension with commutative Galois group) of the field
$\mathbb{Q}$ of rational numbers is contained in a field of roots of
unity. Thus, so to say, certain values of the exponential function
generate the maximal abelian extension of $\mathbb{Q}$. Such an
``explicit'' construction is also possible for an imaginary quadratic
field. One has to use the class invariants and also the values of a
certain function, related to the Weierstrass $\wp$-function
(see §\ref{sec:I-ramified}). This theory is essentially what is called
``complex multiplication''. We shall first deal with unramified
extensions. The results pertaining to this case may be embodied in the
following three theorems, where $K$ is an imaginary quadratic field,
$\mathfrak{k}_{i}$, $(1\leq i\leq h)$, its ideal classes.

\setcounter{theorem}{0}
\renewcommand{\thetheorem}{\Roman{theorem}}

\begin{theorem}\label{thm:I-one}
The class invariants $j(\mathfrak{k}_{i})$ are algebraic integers.
\end{theorem}

\noindent[Let us remark in passing that there is a sort of converse to
Theorem~\ref{thm:I-one}. Namely, C.~L. Siegel
(\emph{Transcendental numbers}, Annals of Math.\ Studies $16$,
Princeton, 1949, pp.~98--99) has deduced from certain results of
Schneider that if $z$ is an algebraic number in the upper half-plane not
belonging to a quadratic imaginary field, then $j(z)$ is transcendental.]

\begin{theorem}\label{thm:I-two}
$K(j(\mathfrak{k}_{i}))$ is independent of $i$, $(1\leq i\leq h)$, and is
the maximal unramified abelian extension of $K$.
\end{theorem}

\noindent(Unramified means that every prime ideal of $K$ decomposes in a
product of distinct prime ideals with exponents $1$.)

By class field theory, it is known that the maximal unramified abelian
extension of $K$ (Hilbert's ``absolute class field'') exists and that its
Galois group $G_{K}$ is canonically isomorphic to the group $C_{K}$ of
ideal classes; the next theorem describes how it operates on the
$j(\mathfrak{k}_{i})$.

\begin{theorem}\label{thm:I-three}
Let $\mathfrak{k}\in C_{K}$ and let $\sigma_{\mathfrak{k}}\in G_{K}$ be
its image by the isomorphism of class field theory. Then
\[
\sigma_{\mathfrak{k}}(j(\mathfrak{k}_{i}))
   =j(\mathfrak{k}^{-1}\cdot\mathfrak{k}_{i})\,.
\]
\end{theorem}

For the proofs of the above theorems, we shall follow the method of the
first part of Hasse's paper quoted in §5, which presupposes class-field
theory. More specifically, Hasse establishes by function theoretical
arguments Theorem~\ref{thm:I-one} and the congruence
\[
j(\mathfrak{p}^{-1}\cdot\mathfrak{k})
   \equiv j(\mathfrak{k})^{N\mathfrak{p}}\mod \mathfrak{p}
\tag{1}\label{eq:I-1}
\]
($\mathfrak{p}$ prime ideal of $K$, $\mathfrak{k}$ ideal class,
$N\mathfrak{p}$ absolute norm of $\mathfrak{p}$), for almost all
(i.e.\ all but a finite number of) prime ideals of $K$ having absolute
degree one (i.e.\ $N\mathfrak{p}=p$ with $p$ prime) with respect to
$\mathbb{Q}$. The validity of \eqref{eq:I-1} for all $\mathfrak{p}$ and
Theorems~\ref{thm:I-two}, \ref{thm:I-three} will then follow by class
field theory.

% ---------------------------------------------------------------
\section{Algebraic interpretation}

For any complex number $j$, there is an elliptic curve $X$ defined over
$\mathbb{Q}(j)$, with invariant $j$; it may for instance be given by the
equation
\[
y^{2}=4x^{3}-h(x+1),\qquad\bigl(h=27j\,(j-2^{6}\!\cdot 3^{3})^{-1}\bigr).
\]

If in particular $j=j(\mathfrak{k})$, with $K,\mathfrak{k}$ as before,
then $K(j)$ is the smallest field of definition for all elements of
$A(X)$. From this, it follows easily that the numbers $j(\mathfrak{k}_{i})$
are algebraic and that $L=K(j(\mathfrak{k}_{i}))$ is independent from
$i$.

Moreover, given a prime ideal $\mathfrak{p}$ of $K$, there is an
algebraic procedure to obtain from $X$ a curve, call it $X_{\mathfrak{p}}$,
with invariant $j(\mathfrak{p}^{-1}\cdot\mathfrak{k})$, also defined over
$K(j)$.

Deuring's method relies on the notion of reduction of $X$ modulo a prime
ideal $\mathfrak{q}$ of $L$. This means that we consider the curve with
the same equation as $X$, but the coefficients being reduced mod
$\mathfrak{q}$. It may be shown that for almost all $\mathfrak{q}$, this
is possible and leads to an elliptic curve $\bar{X}$ which does not depend
on the particular equation chosen to define $X$. The curve $\bar{X}$ is
then defined over a field of characteristic $p$, where $p$ is the prime
number contained in $\mathfrak{q}$.

To any algebraic variety $V$ defined over a field of characteristic
$p\neq 0$, one can associate its ``$p$-th power'' $V^{p}$. Roughly
speaking, $V^{p}$ is obtained by raising to the $p$-th power all
coordinates of the points in any affine model of $V$. In particular we can
consider $\bar{X}^{\mathfrak{q}}$ for any $q=p^{s}$ ($s$ positive integer).

With these notations, Deuring proves by a simple argument that the equality
\[
\bar{x}^{N\mathfrak{p}}=\overline{(x_{\mathfrak{p}})\,},\qquad
(\mathfrak{p}=K\cap\mathfrak{q}),
\tag{3}\label{eq:I-3}
\]
holds for almost all $\mathfrak{p}$ of degree $1$, whence the congruence
\eqref{eq:I-1}.

So far, it has not been possible to prove along these lines that the
class invariants are integral numbers. Deuring's proof for this fact is
the analogue with formal power series of the classical argument.

% ---------------------------------------------------------------
\section{Ramified extensions}\label{sec:I-ramified}

Let as before $K=\mathbb{Q}(\sqrt{-d})$ ($d$ positive square-free
integer) be a quadratic imaginary field, and $X$ an elliptic curve with
invariant $j=j(\mathfrak{k})$, where $\mathfrak{k}$ is an ideal class of
$K$. Its group of automorphisms $\operatorname{Aut} X$ is the group of
units of $A(X)$, i.e.\ of $\mathfrak{o}(K)$. This group is cyclic of
order $2$ (resp.\ $4$, resp.\ $6$) when $d\neq 1,3$ (resp.\ $d=1$,
resp.\ $d=3$).

The quotient $X'$ of $X$ by $\operatorname{Aut} X$ is an algebraic curve
(of genus zero), defined over $K(j)$.

Since $X$ is a group variety, defined over $K(j)$, its points of order
$n$ will be algebraic and so will be their images in $X'$. We can now
state the main theorem of complex multiplication for ramified extensions:

\begin{theorem}\label{thm:I-four}
Let $K$ be an imaginary quadratic field, $j$ a class invariant of $K$,
and $X$ an elliptic curve with modular invariant $j$ defined over
$L=K(j)$. The maximal abelian extension of $K$ can then be obtained by
adjoining to $L$ the coordinates in $X'=X/\operatorname{Aut}X$ of all
points of order $n$ of $X$ $(n=1,2,\ldots)$.
\end{theorem}

In order to express Theorem~\ref{thm:I-four} analytically, we introduce
the function
\[
\tau(u;\omega_{1},\omega_{2})
=(-1)^{e/2}\,\wp^{e/2}(u;\omega_{1},\omega_{2})\,
g^{(e)}(\omega_{1},\omega_{2})
\]
where $e=\text{order }\operatorname{Aut}X$, $\wp(u;\omega_{1},\omega_{2})$
is the Weierstrass function,
\[
\begin{gathered}
g^{(2)}=2^{7}\cdot 3^{5}\,g_{2}\,g_{3}/\Delta\\
g^{(3)}=2^{8}\cdot 3^{4}\,g_{2}^{2}/\Delta\\
g^{(6)}=2^{9}\cdot 3^{6}\,g_{3}/\Delta\,.
\end{gathered}
\]

We have then:
\begin{namedstmt}{THEOREM IV$'$}
Let $K$ be an imaginary quadratic field, $j$ a class invariant of $K$,
$\omega_{1}$, $\omega_{2}$ a basis for some ideal of $K$. Then the maximal
abelian extension of $K$ may be obtained by adjoining to $L=K(j)$ the
numbers
$\tau\bigl((a\omega_{1}+b\omega_{2})/n;\omega_{1},\omega_{2}\bigr)$,
$(a,b,n\in\mathbb{Z},\ n>0)$.
\end{namedstmt}

The maximal abelian extension of $K$ can also be obtained by adjoining
to $K$ the roots of unity, the values $j(z)$ of the modular function for
all $z\in K^{*}$ having positive imaginary part, and square roots of
elements in the field thus obtained. For this and its relation to the
so-called Kronecker Jugendtraum, see Hasse's \emph{Klassenk\"orper
Bericht}
\cite{hasse1926bericht}, §10.

% ---------------------------------------------------------------
\section{Bibliographical notes}\label{sec:I-notes}

We content ourselves with some brief indications, without making any
attempt towards completeness. As to the 19-th century literature, we just
quote \cite{kroneckerWerke}.

The first two systematic and detailed accounts are to be found in
\cite{weber1908algebra,fueter1924moduln,fueter1927moduln}.

Weber's book contains most of the essential results. However, although
Weber already introduces and obtains several properties of the function
$\tau$, both he and Fueter have to use other, more complicated, functions
to generate the maximal abelian extension. That this could be performed
by means of $\tau$ only was first shown by Hasse
\cite{hasse1927neue,hasse1931neue}.

Teil II and the above mentioned books follow the analytical method.
Hasse's Teil I combines analysis and class field theory.

The purely algebraic approach was initiated and carried out by
M.~Deuring. See notably \cite{deuring1947algebraische,deuring1952struktur}.

It was clear from the outset that a main obstacle to a generalization of
Deuring's methods to higher dimensional abelian varieties was the lack of
a good theory for reduction mod $\mathfrak{p}$ in algebraic geometry.
This was recently supplied by Shimura \cite{shimura1955reduction}
(Amer.\ Jour.\ Math.\ 77, 134--176 (1955)), and was applied to higher
dimensional extensions of complex multiplication by Shimura, Taniyama,
Weil. See \cite{shimura1955tokyo,taniyama1955tokyo,weil1955tokyo}.

The two-dimensional case had been considered long ago by Hecke (following
a suggestion of Hilbert), using analytical methods. See
\cite{hecke1912hohere,hecke1913konstruktion}.

[(Added in 1965). The Shimura--Taniyama--Weil theory has been published:
\cite{shimuramataniyama1961}. A systematic exposition of the analytic
method is given in \cite{deuringklassenkorper}. For further results, see
\cite{ramachandra1964applications}.]
